\section{Synthesis: The World as Structure}
\label{sec:synthesis}

\subsection{The Flow of This Theory}

This work began from a single axiom:
\begin{equation}
  \boxed{\text{There is difference.}}
\end{equation}
In the stated representations, difference is associated with gradients,
gradients with an effective flow relation, flow with defined vorticity, and
vorticity with a candidate enclosure mechanism.
\begin{equation}
  \boxed{\text{The model proposes a route from dynamics to structure.}}
\end{equation}

\subsection{The Overall Structure}

The following historical compact display summarizes the program's topic order:
\begin{equation}
  \boxed{
    \text{Difference}
    \;\to\; C_{ij}
    \;\to\; J
    \;\to\; \omega
    \;\to\; L
    \;\to\; g
    \;\to\; \text{Einstein equation}
    \;\to\; \text{Universe}
  }
\end{equation}
\begin{center}
\begin{tabular}{ll}
\toprule
Quantity & Meaning \\
\midrule
$C_{ij}$ & causal log \\
$J$ & flow of logs \\
$\omega$ & vorticity \\
$L$ & closure density \\
$g$ & effective metric \\
Einstein equation & fixed-point structure \\
Universe & global $\sigma > 0$ phase \\
\bottomrule
\end{tabular}
\end{center}
\begin{equation}
  \boxed{\text{Everything forms a single structural chain.}}
\end{equation}

Here ``single structural chain'' means a connected research program, not one
repeated logical relation.  The edge roles are:
\begin{description}
  \item[$C_{ij}\to J_{ij}$] discrete definition;
  \item[$J_{ij}\to\mathbf J(x)$] informal discrete-to-continuum bridge, OPEN;
  \item[$\mathbf J(x)\to\omega$] continuum curl definition;
  \item[$\omega\to L$] effective proportional relation;
  \item[$L,L_{ab}\to L_{\mu\nu}$] carrier / coarse bridge, OPEN;
  \item[network / closure quantities $\to g_{\mu\nu}$] partial effective
        endpoint, with reconstruction and physical interpretation OPEN;
  \item[$g_{\mu\nu}\to$ Einstein equation] intended weak-field structural
        comparison, not exact recovery;
  \item[Einstein comparison $\to$ Universe] proposed cosmological
        interpretation, not an established physical derivation.
\end{description}
Thus the compact arrows do not constitute a new arrow calculus or establish
physical emergence by themselves.

\subsection{Correspondence Across Volumes}

The three volumes form a connected program whose transitions retain the
distinct roles listed above.

\textbf{Vol.~1 --- Generation:}
causal log $C_{ij}$; time $d\tau/dt = \rho$;
space $d_{ij} = -\log\tilde{S}_{ij}$.
Formal time-like and spatial constructions are supplied; their physical
interpretations remain proposed.

\textbf{Vol.~2 --- Structure:}
triangular closure;
gauge symmetry $\mathrm{SU}(3)\times\mathrm{SU}(2)\times\mathrm{U}(1)$;
particles, charge, mass.
Structure stabilizes.

\textbf{Vol.~3 --- Dynamics:}
closure quantities enter an effective gravity-like proposal; the metric
fixed-point relation is a formal target compared structurally with the
weak-field Einstein equation; and the $\sigma > 0$ phase is given a proposed
cosmological interpretation.  The carrier, continuum, exact-recovery, and
observational bridges remain open.
\begin{equation}
  \boxed{\text{Generation, structure, and dynamics are compared in one program.}}
\end{equation}

\subsection{Reinterpreting Conservation Laws}

Conventional physics treats conservation laws as fundamental;
violations appear as exceptions.
Within VED, the continuity equation is:
\begin{equation}
  \partial_t \rho + \nabla \cdot J = \sigma.
\end{equation}
Two regimes appear:
$\sigma = 0$ implies conservation holds;
$\sigma>0$ denotes a positive source regime in this continuum model; an
equivalence with physical generation is not established.
\begin{equation}
  \boxed{\text{The source-free balance relation is a model-phase regime.}}
\end{equation}
This regime classification does not establish whether physical conservation
laws are fundamental or emergent.

\subsection{Reinterpreting Gravity}

Conventional interpretations treat gravity as force or as curvature.
VED interprets it differently.
\begin{equation}
  \boxed{\text{Modified distance is a candidate gravity-like description.}}
\end{equation}
\begin{equation}
  \boxed{\text{Closure density is a candidate gravity source.}}
\end{equation}

\subsection{Black Holes and the Universe}

The black-hole-interior and expanding-universe candidate descriptions both
use a $\sigma>0$ regime.
\begin{equation}
  \boxed{\text{Black-hole interiors and cosmic expansion are proposed as
  analogous } \sigma>0 \text{ regimes.}}
\end{equation}

\subsection{The Final Inversion}

Conventional physics begins by assuming spacetime, fields, and symmetry.
VED instead proposes structural constructions and interpretation maps.
\begin{equation}
  \boxed{\text{Closure structure supplies the common model input.}}
\end{equation}
Time-like ordering and network distance are formally constructed.
Physical time and space remain proposed interpretations.
Particles are represented by candidate closure minima.
Forces are flow.
The universe is a phase.

\subsection{The Final Statement}

\begin{equation}
  \boxed{\text{The world is not closed.}}
\end{equation}
\begin{equation}
  \boxed{\text{It exists by never fully closing.}}
\end{equation}
This statement serves both as the starting point and the conclusion of VED.
\begin{equation}
  \text{Difference}
  \;\to\; C_{ij}
  \;\to\; J
  \;\to\; \omega
  \;\to\; L
  \;\to\; g_{\mu\nu}
  \;\to\; \text{Einstein equation}
  \;\to\; \text{Universe}
\end{equation}
\begin{equation}
  \boxed{\begin{aligned}
    \text{Time-like quantity} &= \text{log accumulation} \\
    \text{Spatial candidate}  &= \text{network log-distance} \\
    \text{Particle-like candidate} &= \text{closure minimum} \\
    \text{Gravity-like candidate}  &= \text{effective closure-density potential} \\
    \text{Cosmological candidate} &= \text{positive-source phase model}
  \end{aligned}}
\end{equation}
\begin{equation}
  \boxed{\text{The model constructions share } C_{ij} \text{ as an upstream input.}}
\end{equation}
